% !TITLE 咏怀古迹·其一
% !AUTHOR 杜甫
% !DYNASTY 唐
% !GENRE 七言律诗
% !ORDER 34

\begin{poem}
支离东北风尘际\textsuperscript{1}，漂泊西南天地间。
三峡楼台淹日月\textsuperscript{2}，五溪\textsuperscript{3}衣服共云山。
羯胡\textsuperscript{4}事主终无赖，词客哀时且未还。
庾信\textsuperscript{5}平生最萧瑟\textsuperscript{6}，暮年诗赋动江关\textsuperscript{7}。
\end{poem}

\begin{annotations}
\notetext{1}{支离东北风尘际：在东北的战乱（安史之乱）中流离失所。支离，破碎、离散。}
\notetext{2}{淹日月：岁月在这里淹留（被拖延）。三峡的幽深遮蔽了日月，暗示时间漫长和被困的处境。}
\notetext{3}{五溪：湘西一带的少数民族聚居地。当地人的衣服五色斑斓，与云山相映。}
\notetext{4}{羯胡：指安禄山等叛军。羯（jié），北方少数民族之一。安禄山是胡人出身。“终无赖”——终究不可信赖。}
\notetext{5}{庾信：南北朝文学家，出使西魏时被扣留，后北周灭梁，再也回不了南方，晚年写下著名的《哀江南赋》。}
\notetext{6}{萧瑟：凄凉、寂寞。庾信一生最是萧瑟——杜甫以此为自况。}
\notetext{7}{动江关：感动了江关（泛指天下）。庾信晚年诗赋的感染力之大，杜甫以此暗喻自己的创作抱负。}
\end{annotations}

\begin{appreciation}
这首诗是杜甫《咏怀古迹五首》中的第一首，写于大历元年（766年），杜甫流寓夔州之时。五首诗分别咏怀庾信、宋玉、王昭君、刘备、诸葛亮，而这一首，是诗人以庾信自况。

支离东北风尘际，漂泊西南天地间——安史之乱中，我在东北风尘中流离失所；如今又漂泊到西南的天地之间。支离是破碎、离散的意思，形容战乱中家园破碎、身心俱疲的状态。杜甫从洛阳到长安，从秦州到成都，再到夔州，一生漂泊，恰如庾信从南朝入北朝，永不得归。

三峡楼台淹日月，五溪衣服共云山——三峡的楼台遮蔽了日月，五溪的蛮夷与云山共处。淹日月，写出了三峡的幽深高峻，也暗示了时间的漫长和诗人被囚禁般的处境。五溪指湘西一带的少数民族地区，衣服共云山，写当地风俗与中原不同，衣饰斑斓，与云山相映。杜甫身处异乡，看什么都觉得陌生。

羯胡事主终无赖，词客哀时且未还——羯胡侍奉君主终究没有好结果，词客哀叹时局却至今不能归还。羯胡指安禄山等叛军，无赖即不可信赖。这里杜甫在骂叛军，也在感慨自己：因为战乱，因为时代的动荡，他这样的文人只能漂泊在外，有家难回。

庾信平生最萧瑟，暮年诗赋动江关——庾信一生最为萧瑟凄凉，到了暮年，他的诗赋却感动了江关。庾信是南朝梁代的文学家，出使西魏时被扣留，后来北周灭梁，他再也回不去南方，只能在北方终老。他晚年写下了著名的《哀江南赋》，哀悼故国，字字血泪。萧瑟是凄凉、寂寞的意思。动江关，是说他的诗赋感人至深，连江关都为之震动。

杜甫与庾信，相隔三百年，命运却惊人地相似：都经历了国家的破碎，都漂泊异乡，都在晚年写出了最动人的诗篇。《哀江南赋》中，庾信写道：日暮途远，人间何世。杜甫读到这些文字，大概会觉得是在写自己。他们都是时代的弃儿，却也是时代的见证者。最萧瑟的平生，换来了最动人的诗赋——这不知是幸运，还是不幸。
\end{appreciation}
